% ======================================================
% RESULTADOS PARCIALES
% Se incluye desde main.tex con % ======================================================
% RESULTADOS PARCIALES
% Se incluye desde main.tex con % ======================================================
% RESULTADOS PARCIALES
% Se incluye desde main.tex con % ======================================================
% RESULTADOS PARCIALES
% Se incluye desde main.tex con \input{resultados_parciales}, después de Costos
% y antes de las Referencias.
% ======================================================
\providecommand{\completar}[1]{\textcolor{red}{\textbf{[#1]}}}
\providecommand{\chipRP}[2]{\tikz[baseline=(c.base)]\node[fill=#1, text=white,
    rounded corners=2pt, inner sep=2.5pt, font=\scriptsize\bfseries\sffamily] (c) {#2};}
\providecommand{\estCUMPLIDO}{\chipRP{green!45!black}{CUMPLIDO}}
\providecommand{\estPARCIAL}{\chipRP{blue!65!black}{PARCIAL}}
\providecommand{\estPENDIENTE}{\chipRP{orange!85!black}{PENDIENTE}}

\section{Resultados parciales}
\label{sec:resultados}

En el primer parcial se fijaron cinco objetivos específicos, uno por área, cada
uno con un criterio de aceptación medible. La tabla resume el resultado
obtenido al 23/11/2026 y lo que queda para el informe final.

\begin{table}[H]
    \centering
    \small
    \caption{Objetivos específicos: criterio de aceptación y resultado al segundo parcial.}
    \label{tab:resultados_parciales}
    \begin{tabularx}{\textwidth}{@{}>{\raggedright\arraybackslash}p{1.9cm} >{\raggedright\arraybackslash}p{3.6cm} L c@{}}
        \toprule
        \textbf{Área} & \textbf{Criterio de aceptación} & \textbf{Resultado obtenido} & \textbf{Estado} \\
        \midrule
        Electrónica &
        Consumo medido por módulo y autonomía medida con la UPS. &
        Radar: 126~mA a 12~V ($\approx$1,5~W), sin modo de reposo.
        Autonomía medida: 2~h~30~min, unas 8 veces menos que lo estimado, por
        las baterías de la UPS (sección~\ref{sec:autonomia}). Falta medir el
        ESP32 y el módem. &
        \estPARCIAL \\
        \addlinespace
        Firmware &
        0 perdidos y 0 duplicados en $\geq$10 cortes; disponibilidad
        $\geq$95\,\% en 7 días. &
        Corte de red n.º~1: 7 de 7 recuperadas, 0 perdidas, 0 duplicadas.
        Corte de energía n.º~1 (fin de la batería): el id siguió sin saltos
        ni repeticiones. En 65~h: 775 de 786 turnos medidos (98,6\,\%; los
        11 que faltan son la UPS agotada), 770 en la página, 0 repetidos.
        Faltan 9 cortes de red, 9 de energía y completar los 7 días
        (27/11). &
        \estPARCIAL \\
        \addlinespace
        Montaje &
        Antena nivelada y haz libre de obstáculos. &
        Radar montado con su kit sobre un perfil en voladizo; nivelado con
        la burbuja en los dos ejes, centrado sobre el balde y a menos de
        1,2~m del agua. &
        \estCUMPLIDO \\
        \addlinespace
        Interfaz &
        Web con históricos, antigüedad del dato, estado del sensor y alarma. &
        Página publicada en GitHub Pages con dato actual, antigüedad,
        estado EN LÍNEA/CAÍDA, estado del sensor, alarma e históricos de
        24~h, 7 y 30 días (figura~\ref{fig:web}). El estado de falla se
        observó el 20/11 con el radar desconectado. &
        \estCUMPLIDO \\
        \addlinespace
        Validación &
        Error $\leq\pm$10~mm frente a cinta métrica. &
        Un punto de control sobre superficie sólida: 0,197~m del radar
        contra 19,7~cm con cinta. Falta el ensayo sobre agua. &
        \estPENDIENTE \\
        \bottomrule
    \end{tabularx}
\end{table}

\subsection*{Comparación con la línea de base}

La estación de Proyecto~3 tuvo una disponibilidad de 88,2\,\% en nueve meses
y no permitía saber si un hueco se debía al sensor, a la red o a la energía.
En 65 horas de funcionamiento continuo con el radar la estación midió el
98,6\,\% de los turnos, y el 100\,\% de los turnos en que tuvo energía; gracias
al id correlativo, cada dato faltante se pudo explicar: o la estación estaba
apagada (UPS agotada) o la lectura se midió y no llegó a la página. El período
todavía es corto para compararlo con los nueve meses de Proyecto~3; por eso se
completa la prueba de 7 días.

\subsection*{Ensayos reprogramados}

En la respuesta a la evaluación del primer parcial se previó adelantar los
cortes, la disponibilidad y la autonomía con el sensor de Proyecto~3 mientras
llegaba el radar. Esos ensayos no se hicieron con el sensor anterior: se
empezaron con la estación ya integrada con el radar, y al cierre de este
parcial están hechos el primer corte de red, el primer corte de energía y el
ensayo de autonomía. La campaña comprometida no está terminada; se mantienen
los mismos criterios de aceptación y su verificación se completa para el
informe final, con los ensayos de la tabla siguiente.

\subsection*{Pendiente para el informe final}

\begin{table}[H]
    \centering
    \small
    \caption{Ensayos que faltan y cómo se harán.}
    \begin{tabularx}{\textwidth}{@{}>{\raggedright\arraybackslash}p{3cm} L >{\raggedright\arraybackslash}p{2.6cm}@{}}
        \toprule
        \textbf{Ensayo} & \textbf{Cómo} & \textbf{Umbral} \\
        \midrule
        Exactitud sobre agua & Modo banco, 4 alturas $\times$ 10 lecturas
            sobre el balde, frente a cinta métrica. & $\leq\pm$10~mm \\
        Cortes de red & 9 cortes más, incluidos módem sin señal y WiFi caído
            con el módem enviando. & 0 perdidos, 0 duplicados \\
        Cortes de energía & 9 cortes más, incluso durante la escritura en la
            microSD y el envío. & 0 perdidos, 0 duplicados \\
        Disponibilidad & 7 días continuos sin intervenir la estación. &
            $\geq$95\,\% \\
        Consumo y autonomía & Corriente del ESP32 y del módem (media y
            picos al transmitir); repetir la descarga con baterías nuevas o
            una batería mayor. & Se informa \\
        \bottomrule
    \end{tabularx}
\end{table}

\subsection*{Cronograma hasta el informe final}

La estación sigue funcionando en el laboratorio. Los ensayos que interrumpen la
estación (cortes, reprogramar en modo banco, cambiar las baterías) se hacen
después de cerrar los 7 días continuos, para no mezclar sus efectos.

\begin{table}[H]
    \centering
    \small
    \caption{Cronograma hasta la entrega del informe final (14/12/2026).}
    \begin{tabularx}{\textwidth}{@{}>{\raggedright\arraybackslash}p{2.6cm}L@{}}
        \toprule
        \textbf{Semana} & \textbf{Actividades} \\
        \midrule
        24/11 -- 30/11 & Cerrar la prueba de 7 días sin intervenir la estación
            (27/11 22:00) y conciliar microSD, base de datos y página. Desde el
            28/11: ensayo de exactitud sobre agua (4 alturas $\times$ 10
            lecturas). \\
        1/12 -- 7/12 & 9 cortes de red (incluidos módem sin señal y WiFi caído
            con el módem enviando) y 9 de energía, también durante la escritura
            en la microSD. Consumo del ESP32 y del módem. Cambio de las
            baterías de la UPS y nuevo ensayo de autonomía. \\
        8/12 -- 14/12 & Análisis de resultados frente a los criterios de
            aceptación y redacción del informe final. Entrega: 14/12. \\
        \bottomrule
    \end{tabularx}
\end{table}

\subsection*{Anexos}

Las fotos de los ensayos, de los circuitos, del montaje mecánico, de la
configuración del radar, de la página de monitoreo y de la base de datos se
reúnen en una carpeta del repositorio, ordenadas por tema.

\begin{itemize}
    \item Fotos (anexo):
        \video{https://github.com/HecVelaz/Sistema-de-adquisicion-y-monitoreo-de-nivel-P4/tree/main/documentos/Reportes/P4_2doParcial/Fotos}{carpeta
        Fotos del repositorio}
    \item Videos de las pruebas:
        \video{https://www.youtube.com/playlist?list=PLdK9Oc0zBIqY}{lista de
        reproducción en YouTube}
    \item Página de monitoreo:
        \video{https://hecvelaz.github.io/Sistema-de-adquisicion-y-monitoreo-de-nivel-P4/}{Estación
        Mburicaó (GitHub Pages)}
    \item Repositorio del proyecto:
        \video{https://github.com/HecVelaz/Sistema-de-adquisicion-y-monitoreo-de-nivel-P4}{repositorio
        en GitHub}
\end{itemize}
, después de Costos
% y antes de las Referencias.
% ======================================================
\providecommand{\completar}[1]{\textcolor{red}{\textbf{[#1]}}}
\providecommand{\chipRP}[2]{\tikz[baseline=(c.base)]\node[fill=#1, text=white,
    rounded corners=2pt, inner sep=2.5pt, font=\scriptsize\bfseries\sffamily] (c) {#2};}
\providecommand{\estCUMPLIDO}{\chipRP{green!45!black}{CUMPLIDO}}
\providecommand{\estPARCIAL}{\chipRP{blue!65!black}{PARCIAL}}
\providecommand{\estPENDIENTE}{\chipRP{orange!85!black}{PENDIENTE}}

\section{Resultados parciales}
\label{sec:resultados}

En el primer parcial se fijaron cinco objetivos específicos, uno por área, cada
uno con un criterio de aceptación medible. La tabla resume el resultado
obtenido al 23/11/2026 y lo que queda para el informe final.

\begin{table}[H]
    \centering
    \small
    \caption{Objetivos específicos: criterio de aceptación y resultado al segundo parcial.}
    \label{tab:resultados_parciales}
    \begin{tabularx}{\textwidth}{@{}>{\raggedright\arraybackslash}p{1.9cm} >{\raggedright\arraybackslash}p{3.6cm} L c@{}}
        \toprule
        \textbf{Área} & \textbf{Criterio de aceptación} & \textbf{Resultado obtenido} & \textbf{Estado} \\
        \midrule
        Electrónica &
        Consumo medido por módulo y autonomía medida con la UPS. &
        Radar: 126~mA a 12~V ($\approx$1,5~W), sin modo de reposo.
        Autonomía medida: 2~h~30~min, unas 8 veces menos que lo estimado, por
        las baterías de la UPS (sección~\ref{sec:autonomia}). Falta medir el
        ESP32 y el módem. &
        \estPARCIAL \\
        \addlinespace
        Firmware &
        0 perdidos y 0 duplicados en $\geq$10 cortes; disponibilidad
        $\geq$95\,\% en 7 días. &
        Corte de red n.º~1: 7 de 7 recuperadas, 0 perdidas, 0 duplicadas.
        Corte de energía n.º~1 (fin de la batería): el id siguió sin saltos
        ni repeticiones. En 65~h: 775 de 786 turnos medidos (98,6\,\%; los
        11 que faltan son la UPS agotada), 770 en la página, 0 repetidos.
        Faltan 9 cortes de red, 9 de energía y completar los 7 días
        (27/11). &
        \estPARCIAL \\
        \addlinespace
        Montaje &
        Antena nivelada y haz libre de obstáculos. &
        Radar montado con su kit sobre un perfil en voladizo; nivelado con
        la burbuja en los dos ejes, centrado sobre el balde y a menos de
        1,2~m del agua. &
        \estCUMPLIDO \\
        \addlinespace
        Interfaz &
        Web con históricos, antigüedad del dato, estado del sensor y alarma. &
        Página publicada en GitHub Pages con dato actual, antigüedad,
        estado EN LÍNEA/CAÍDA, estado del sensor, alarma e históricos de
        24~h, 7 y 30 días (figura~\ref{fig:web}). El estado de falla se
        observó el 20/11 con el radar desconectado. &
        \estCUMPLIDO \\
        \addlinespace
        Validación &
        Error $\leq\pm$10~mm frente a cinta métrica. &
        Un punto de control sobre superficie sólida: 0,197~m del radar
        contra 19,7~cm con cinta. Falta el ensayo sobre agua. &
        \estPENDIENTE \\
        \bottomrule
    \end{tabularx}
\end{table}

\subsection*{Comparación con la línea de base}

La estación de Proyecto~3 tuvo una disponibilidad de 88,2\,\% en nueve meses
y no permitía saber si un hueco se debía al sensor, a la red o a la energía.
En 65 horas de funcionamiento continuo con el radar la estación midió el
98,6\,\% de los turnos, y el 100\,\% de los turnos en que tuvo energía; gracias
al id correlativo, cada dato faltante se pudo explicar: o la estación estaba
apagada (UPS agotada) o la lectura se midió y no llegó a la página. El período
todavía es corto para compararlo con los nueve meses de Proyecto~3; por eso se
completa la prueba de 7 días.

\subsection*{Ensayos reprogramados}

En la respuesta a la evaluación del primer parcial se previó adelantar los
cortes, la disponibilidad y la autonomía con el sensor de Proyecto~3 mientras
llegaba el radar. Esos ensayos no se hicieron con el sensor anterior: se
empezaron con la estación ya integrada con el radar, y al cierre de este
parcial están hechos el primer corte de red, el primer corte de energía y el
ensayo de autonomía. La campaña comprometida no está terminada; se mantienen
los mismos criterios de aceptación y su verificación se completa para el
informe final, con los ensayos de la tabla siguiente.

\subsection*{Pendiente para el informe final}

\begin{table}[H]
    \centering
    \small
    \caption{Ensayos que faltan y cómo se harán.}
    \begin{tabularx}{\textwidth}{@{}>{\raggedright\arraybackslash}p{3cm} L >{\raggedright\arraybackslash}p{2.6cm}@{}}
        \toprule
        \textbf{Ensayo} & \textbf{Cómo} & \textbf{Umbral} \\
        \midrule
        Exactitud sobre agua & Modo banco, 4 alturas $\times$ 10 lecturas
            sobre el balde, frente a cinta métrica. & $\leq\pm$10~mm \\
        Cortes de red & 9 cortes más, incluidos módem sin señal y WiFi caído
            con el módem enviando. & 0 perdidos, 0 duplicados \\
        Cortes de energía & 9 cortes más, incluso durante la escritura en la
            microSD y el envío. & 0 perdidos, 0 duplicados \\
        Disponibilidad & 7 días continuos sin intervenir la estación. &
            $\geq$95\,\% \\
        Consumo y autonomía & Corriente del ESP32 y del módem (media y
            picos al transmitir); repetir la descarga con baterías nuevas o
            una batería mayor. & Se informa \\
        \bottomrule
    \end{tabularx}
\end{table}

\subsection*{Cronograma hasta el informe final}

La estación sigue funcionando en el laboratorio. Los ensayos que interrumpen la
estación (cortes, reprogramar en modo banco, cambiar las baterías) se hacen
después de cerrar los 7 días continuos, para no mezclar sus efectos.

\begin{table}[H]
    \centering
    \small
    \caption{Cronograma hasta la entrega del informe final (14/12/2026).}
    \begin{tabularx}{\textwidth}{@{}>{\raggedright\arraybackslash}p{2.6cm}L@{}}
        \toprule
        \textbf{Semana} & \textbf{Actividades} \\
        \midrule
        24/11 -- 30/11 & Cerrar la prueba de 7 días sin intervenir la estación
            (27/11 22:00) y conciliar microSD, base de datos y página. Desde el
            28/11: ensayo de exactitud sobre agua (4 alturas $\times$ 10
            lecturas). \\
        1/12 -- 7/12 & 9 cortes de red (incluidos módem sin señal y WiFi caído
            con el módem enviando) y 9 de energía, también durante la escritura
            en la microSD. Consumo del ESP32 y del módem. Cambio de las
            baterías de la UPS y nuevo ensayo de autonomía. \\
        8/12 -- 14/12 & Análisis de resultados frente a los criterios de
            aceptación y redacción del informe final. Entrega: 14/12. \\
        \bottomrule
    \end{tabularx}
\end{table}

\subsection*{Anexos}

Las fotos de los ensayos, de los circuitos, del montaje mecánico, de la
configuración del radar, de la página de monitoreo y de la base de datos se
reúnen en una carpeta del repositorio, ordenadas por tema.

\begin{itemize}
    \item Fotos (anexo):
        \video{https://github.com/HecVelaz/Sistema-de-adquisicion-y-monitoreo-de-nivel-P4/tree/main/documentos/Reportes/P4_2doParcial/Fotos}{carpeta
        Fotos del repositorio}
    \item Videos de las pruebas:
        \video{https://www.youtube.com/playlist?list=PLdK9Oc0zBIqY}{lista de
        reproducción en YouTube}
    \item Página de monitoreo:
        \video{https://hecvelaz.github.io/Sistema-de-adquisicion-y-monitoreo-de-nivel-P4/}{Estación
        Mburicaó (GitHub Pages)}
    \item Repositorio del proyecto:
        \video{https://github.com/HecVelaz/Sistema-de-adquisicion-y-monitoreo-de-nivel-P4}{repositorio
        en GitHub}
\end{itemize}
, después de Costos
% y antes de las Referencias.
% ======================================================
\providecommand{\completar}[1]{\textcolor{red}{\textbf{[#1]}}}
\providecommand{\chipRP}[2]{\tikz[baseline=(c.base)]\node[fill=#1, text=white,
    rounded corners=2pt, inner sep=2.5pt, font=\scriptsize\bfseries\sffamily] (c) {#2};}
\providecommand{\estCUMPLIDO}{\chipRP{green!45!black}{CUMPLIDO}}
\providecommand{\estPARCIAL}{\chipRP{blue!65!black}{PARCIAL}}
\providecommand{\estPENDIENTE}{\chipRP{orange!85!black}{PENDIENTE}}

\section{Resultados parciales}
\label{sec:resultados}

En el primer parcial se fijaron cinco objetivos específicos, uno por área, cada
uno con un criterio de aceptación medible. La tabla resume el resultado
obtenido al 23/11/2026 y lo que queda para el informe final.

\begin{table}[H]
    \centering
    \small
    \caption{Objetivos específicos: criterio de aceptación y resultado al segundo parcial.}
    \label{tab:resultados_parciales}
    \begin{tabularx}{\textwidth}{@{}>{\raggedright\arraybackslash}p{1.9cm} >{\raggedright\arraybackslash}p{3.6cm} L c@{}}
        \toprule
        \textbf{Área} & \textbf{Criterio de aceptación} & \textbf{Resultado obtenido} & \textbf{Estado} \\
        \midrule
        Electrónica &
        Consumo medido por módulo y autonomía medida con la UPS. &
        Radar: 126~mA a 12~V ($\approx$1,5~W), sin modo de reposo.
        Autonomía medida: 2~h~30~min, unas 8 veces menos que lo estimado, por
        las baterías de la UPS (sección~\ref{sec:autonomia}). Falta medir el
        ESP32 y el módem. &
        \estPARCIAL \\
        \addlinespace
        Firmware &
        0 perdidos y 0 duplicados en $\geq$10 cortes; disponibilidad
        $\geq$95\,\% en 7 días. &
        Corte de red n.º~1: 7 de 7 recuperadas, 0 perdidas, 0 duplicadas.
        Corte de energía n.º~1 (fin de la batería): el id siguió sin saltos
        ni repeticiones. En 65~h: 775 de 786 turnos medidos (98,6\,\%; los
        11 que faltan son la UPS agotada), 770 en la página, 0 repetidos.
        Faltan 9 cortes de red, 9 de energía y completar los 7 días
        (27/11). &
        \estPARCIAL \\
        \addlinespace
        Montaje &
        Antena nivelada y haz libre de obstáculos. &
        Radar montado con su kit sobre un perfil en voladizo; nivelado con
        la burbuja en los dos ejes, centrado sobre el balde y a menos de
        1,2~m del agua. &
        \estCUMPLIDO \\
        \addlinespace
        Interfaz &
        Web con históricos, antigüedad del dato, estado del sensor y alarma. &
        Página publicada en GitHub Pages con dato actual, antigüedad,
        estado EN LÍNEA/CAÍDA, estado del sensor, alarma e históricos de
        24~h, 7 y 30 días (figura~\ref{fig:web}). El estado de falla se
        observó el 20/11 con el radar desconectado. &
        \estCUMPLIDO \\
        \addlinespace
        Validación &
        Error $\leq\pm$10~mm frente a cinta métrica. &
        Un punto de control sobre superficie sólida: 0,197~m del radar
        contra 19,7~cm con cinta. Falta el ensayo sobre agua. &
        \estPENDIENTE \\
        \bottomrule
    \end{tabularx}
\end{table}

\subsection*{Comparación con la línea de base}

La estación de Proyecto~3 tuvo una disponibilidad de 88,2\,\% en nueve meses
y no permitía saber si un hueco se debía al sensor, a la red o a la energía.
En 65 horas de funcionamiento continuo con el radar la estación midió el
98,6\,\% de los turnos, y el 100\,\% de los turnos en que tuvo energía; gracias
al id correlativo, cada dato faltante se pudo explicar: o la estación estaba
apagada (UPS agotada) o la lectura se midió y no llegó a la página. El período
todavía es corto para compararlo con los nueve meses de Proyecto~3; por eso se
completa la prueba de 7 días.

\subsection*{Ensayos reprogramados}

En la respuesta a la evaluación del primer parcial se previó adelantar los
cortes, la disponibilidad y la autonomía con el sensor de Proyecto~3 mientras
llegaba el radar. Esos ensayos no se hicieron con el sensor anterior: se
empezaron con la estación ya integrada con el radar, y al cierre de este
parcial están hechos el primer corte de red, el primer corte de energía y el
ensayo de autonomía. La campaña comprometida no está terminada; se mantienen
los mismos criterios de aceptación y su verificación se completa para el
informe final, con los ensayos de la tabla siguiente.

\subsection*{Pendiente para el informe final}

\begin{table}[H]
    \centering
    \small
    \caption{Ensayos que faltan y cómo se harán.}
    \begin{tabularx}{\textwidth}{@{}>{\raggedright\arraybackslash}p{3cm} L >{\raggedright\arraybackslash}p{2.6cm}@{}}
        \toprule
        \textbf{Ensayo} & \textbf{Cómo} & \textbf{Umbral} \\
        \midrule
        Exactitud sobre agua & Modo banco, 4 alturas $\times$ 10 lecturas
            sobre el balde, frente a cinta métrica. & $\leq\pm$10~mm \\
        Cortes de red & 9 cortes más, incluidos módem sin señal y WiFi caído
            con el módem enviando. & 0 perdidos, 0 duplicados \\
        Cortes de energía & 9 cortes más, incluso durante la escritura en la
            microSD y el envío. & 0 perdidos, 0 duplicados \\
        Disponibilidad & 7 días continuos sin intervenir la estación. &
            $\geq$95\,\% \\
        Consumo y autonomía & Corriente del ESP32 y del módem (media y
            picos al transmitir); repetir la descarga con baterías nuevas o
            una batería mayor. & Se informa \\
        \bottomrule
    \end{tabularx}
\end{table}

\subsection*{Cronograma hasta el informe final}

La estación sigue funcionando en el laboratorio. Los ensayos que interrumpen la
estación (cortes, reprogramar en modo banco, cambiar las baterías) se hacen
después de cerrar los 7 días continuos, para no mezclar sus efectos.

\begin{table}[H]
    \centering
    \small
    \caption{Cronograma hasta la entrega del informe final (14/12/2026).}
    \begin{tabularx}{\textwidth}{@{}>{\raggedright\arraybackslash}p{2.6cm}L@{}}
        \toprule
        \textbf{Semana} & \textbf{Actividades} \\
        \midrule
        24/11 -- 30/11 & Cerrar la prueba de 7 días sin intervenir la estación
            (27/11 22:00) y conciliar microSD, base de datos y página. Desde el
            28/11: ensayo de exactitud sobre agua (4 alturas $\times$ 10
            lecturas). \\
        1/12 -- 7/12 & 9 cortes de red (incluidos módem sin señal y WiFi caído
            con el módem enviando) y 9 de energía, también durante la escritura
            en la microSD. Consumo del ESP32 y del módem. Cambio de las
            baterías de la UPS y nuevo ensayo de autonomía. \\
        8/12 -- 14/12 & Análisis de resultados frente a los criterios de
            aceptación y redacción del informe final. Entrega: 14/12. \\
        \bottomrule
    \end{tabularx}
\end{table}

\subsection*{Anexos}

Las fotos de los ensayos, de los circuitos, del montaje mecánico, de la
configuración del radar, de la página de monitoreo y de la base de datos se
reúnen en una carpeta del repositorio, ordenadas por tema.

\begin{itemize}
    \item Fotos (anexo):
        \video{https://github.com/HecVelaz/Sistema-de-adquisicion-y-monitoreo-de-nivel-P4/tree/main/documentos/Reportes/P4_2doParcial/Fotos}{carpeta
        Fotos del repositorio}
    \item Videos de las pruebas:
        \video{https://www.youtube.com/playlist?list=PLdK9Oc0zBIqY}{lista de
        reproducción en YouTube}
    \item Página de monitoreo:
        \video{https://hecvelaz.github.io/Sistema-de-adquisicion-y-monitoreo-de-nivel-P4/}{Estación
        Mburicaó (GitHub Pages)}
    \item Repositorio del proyecto:
        \video{https://github.com/HecVelaz/Sistema-de-adquisicion-y-monitoreo-de-nivel-P4}{repositorio
        en GitHub}
\end{itemize}
, después de Costos
% y antes de las Referencias.
% ======================================================
\providecommand{\completar}[1]{\textcolor{red}{\textbf{[#1]}}}
\providecommand{\chipRP}[2]{\tikz[baseline=(c.base)]\node[fill=#1, text=white,
    rounded corners=2pt, inner sep=2.5pt, font=\scriptsize\bfseries\sffamily] (c) {#2};}
\providecommand{\estCUMPLIDO}{\chipRP{green!45!black}{CUMPLIDO}}
\providecommand{\estPARCIAL}{\chipRP{blue!65!black}{PARCIAL}}
\providecommand{\estPENDIENTE}{\chipRP{orange!85!black}{PENDIENTE}}

\section{Resultados parciales}
\label{sec:resultados}

En el primer parcial se fijaron cinco objetivos específicos, uno por área, cada
uno con un criterio de aceptación medible. La tabla resume el resultado
obtenido al 23/11/2026 y lo que queda para el informe final.

\begin{table}[H]
    \centering
    \small
    \caption{Objetivos específicos: criterio de aceptación y resultado al segundo parcial.}
    \label{tab:resultados_parciales}
    \begin{tabularx}{\textwidth}{@{}>{\raggedright\arraybackslash}p{1.9cm} >{\raggedright\arraybackslash}p{3.6cm} L c@{}}
        \toprule
        \textbf{Área} & \textbf{Criterio de aceptación} & \textbf{Resultado obtenido} & \textbf{Estado} \\
        \midrule
        Electrónica &
        Consumo medido por módulo y autonomía medida con la UPS. &
        Radar: 126~mA a 12~V ($\approx$1,5~W), sin modo de reposo.
        Autonomía medida: 2~h~30~min, unas 8 veces menos que lo estimado, por
        las baterías de la UPS (sección~\ref{sec:autonomia}). Falta medir el
        ESP32 y el módem. &
        \estPARCIAL \\
        \addlinespace
        Firmware &
        0 perdidos y 0 duplicados en $\geq$10 cortes; disponibilidad
        $\geq$95\,\% en 7 días. &
        Corte de red n.º~1: 7 de 7 recuperadas, 0 perdidas, 0 duplicadas.
        Corte de energía n.º~1 (fin de la batería): el id siguió sin saltos
        ni repeticiones. En 65~h: 775 de 786 turnos medidos (98,6\,\%; los
        11 que faltan son la UPS agotada), 770 en la página, 0 repetidos.
        Faltan 9 cortes de red, 9 de energía y completar los 7 días
        (27/11). &
        \estPARCIAL \\
        \addlinespace
        Montaje &
        Antena nivelada y haz libre de obstáculos. &
        Radar montado con su kit sobre un perfil en voladizo; nivelado con
        la burbuja en los dos ejes, centrado sobre el balde y a menos de
        1,2~m del agua. &
        \estCUMPLIDO \\
        \addlinespace
        Interfaz &
        Web con históricos, antigüedad del dato, estado del sensor y alarma. &
        Página publicada en GitHub Pages con dato actual, antigüedad,
        estado EN LÍNEA/CAÍDA, estado del sensor, alarma e históricos de
        24~h, 7 y 30 días (figura~\ref{fig:web}). El estado de falla se
        observó el 20/11 con el radar desconectado. &
        \estCUMPLIDO \\
        \addlinespace
        Validación &
        Error $\leq\pm$10~mm frente a cinta métrica. &
        Un punto de control sobre superficie sólida: 0,197~m del radar
        contra 19,7~cm con cinta. Falta el ensayo sobre agua. &
        \estPENDIENTE \\
        \bottomrule
    \end{tabularx}
\end{table}

\subsection*{Comparación con la línea de base}

La estación de Proyecto~3 tuvo una disponibilidad de 88,2\,\% en nueve meses
y no permitía saber si un hueco se debía al sensor, a la red o a la energía.
En 65 horas de funcionamiento continuo con el radar la estación midió el
98,6\,\% de los turnos, y el 100\,\% de los turnos en que tuvo energía; gracias
al id correlativo, cada dato faltante se pudo explicar: o la estación estaba
apagada (UPS agotada) o la lectura se midió y no llegó a la página. El período
todavía es corto para compararlo con los nueve meses de Proyecto~3; por eso se
completa la prueba de 7 días.

\subsection*{Ensayos reprogramados}

En la respuesta a la evaluación del primer parcial se previó adelantar los
cortes, la disponibilidad y la autonomía con el sensor de Proyecto~3 mientras
llegaba el radar. Esos ensayos no se hicieron con el sensor anterior: se
empezaron con la estación ya integrada con el radar, y al cierre de este
parcial están hechos el primer corte de red, el primer corte de energía y el
ensayo de autonomía. La campaña comprometida no está terminada; se mantienen
los mismos criterios de aceptación y su verificación se completa para el
informe final, con los ensayos de la tabla siguiente.

\subsection*{Pendiente para el informe final}

\begin{table}[H]
    \centering
    \small
    \caption{Ensayos que faltan y cómo se harán.}
    \begin{tabularx}{\textwidth}{@{}>{\raggedright\arraybackslash}p{3cm} L >{\raggedright\arraybackslash}p{2.6cm}@{}}
        \toprule
        \textbf{Ensayo} & \textbf{Cómo} & \textbf{Umbral} \\
        \midrule
        Exactitud sobre agua & Modo banco, 4 alturas $\times$ 10 lecturas
            sobre el balde, frente a cinta métrica. & $\leq\pm$10~mm \\
        Cortes de red & 9 cortes más, incluidos módem sin señal y WiFi caído
            con el módem enviando. & 0 perdidos, 0 duplicados \\
        Cortes de energía & 9 cortes más, incluso durante la escritura en la
            microSD y el envío. & 0 perdidos, 0 duplicados \\
        Disponibilidad & 7 días continuos sin intervenir la estación. &
            $\geq$95\,\% \\
        Consumo y autonomía & Corriente del ESP32 y del módem (media y
            picos al transmitir); repetir la descarga con baterías nuevas o
            una batería mayor. & Se informa \\
        \bottomrule
    \end{tabularx}
\end{table}

\subsection*{Cronograma hasta el informe final}

La estación sigue funcionando en el laboratorio. Los ensayos que interrumpen la
estación (cortes, reprogramar en modo banco, cambiar las baterías) se hacen
después de cerrar los 7 días continuos, para no mezclar sus efectos.

\begin{table}[H]
    \centering
    \small
    \caption{Cronograma hasta la entrega del informe final (14/12/2026).}
    \begin{tabularx}{\textwidth}{@{}>{\raggedright\arraybackslash}p{2.6cm}L@{}}
        \toprule
        \textbf{Semana} & \textbf{Actividades} \\
        \midrule
        24/11 -- 30/11 & Cerrar la prueba de 7 días sin intervenir la estación
            (27/11 22:00) y conciliar microSD, base de datos y página. Desde el
            28/11: ensayo de exactitud sobre agua (4 alturas $\times$ 10
            lecturas). \\
        1/12 -- 7/12 & 9 cortes de red (incluidos módem sin señal y WiFi caído
            con el módem enviando) y 9 de energía, también durante la escritura
            en la microSD. Consumo del ESP32 y del módem. Cambio de las
            baterías de la UPS y nuevo ensayo de autonomía. \\
        8/12 -- 14/12 & Análisis de resultados frente a los criterios de
            aceptación y redacción del informe final. Entrega: 14/12. \\
        \bottomrule
    \end{tabularx}
\end{table}

\subsection*{Anexos}

Las fotos de los ensayos, de los circuitos, del montaje mecánico, de la
configuración del radar, de la página de monitoreo y de la base de datos se
reúnen en una carpeta del repositorio, ordenadas por tema.

\begin{itemize}
    \item Fotos (anexo):
        \video{https://github.com/HecVelaz/Sistema-de-adquisicion-y-monitoreo-de-nivel-P4/tree/main/documentos/Reportes/P4_2doParcial/Fotos}{carpeta
        Fotos del repositorio}
    \item Videos de las pruebas:
        \video{https://www.youtube.com/playlist?list=PLdK9Oc0zBIqY}{lista de
        reproducción en YouTube}
    \item Página de monitoreo:
        \video{https://hecvelaz.github.io/Sistema-de-adquisicion-y-monitoreo-de-nivel-P4/}{Estación
        Mburicaó (GitHub Pages)}
    \item Repositorio del proyecto:
        \video{https://github.com/HecVelaz/Sistema-de-adquisicion-y-monitoreo-de-nivel-P4}{repositorio
        en GitHub}
\end{itemize}
